\section{The observed terminal class on the complete original window}
This section receives the new continuation's equations (1)--(34).
It verifies its minimum argument in the original coordinates and
strengthens its arithmetic--Gamma same-class comparison. The period
nonvanishing calculation is proved separately in OC below.

Retain a simple hypothetical zeta quartet with
$\rho=1/2+\delta+i\gamma$, $0<\delta<1/2$, and the original fixed
period $\varpi$. Let $k\equiv1\pmod4$, $k\ge9$,
$q=(k+1)^2$, $c=k/2$, and
\[
 \chi_k(S)=\prod_{a,b=0}^k[S-c-(2a-k)\delta-i(2b-k)\gamma],\quad
 v_+(S)=\frac{\chi_k(S)}{(S-k\rho)\chi_k'(k\rho)}.\tag{OB1}
\]
The roots are distinct. The physical source metric is
$dm_k(y)=w_h^{*k}(y)dy$, $S=c+iy$, with $w_h$ in (LR2).
Let $E_k=\mathbb C[S]/(\chi_k)$ and $G_N$ its attained metric at
polynomial degree at most $N$. The observation $\Lambda_k:E_k\to B_k$
is the original period-dependent invariant observation, onto its image.
It includes the original conductor-multiple invariant subspace, as
expressed by FW46. Define
\[
 E_N=v_+^*G_Nv_+,\quad Q_N=(\Lambda_kG_N^{-1}\Lambda_k^*)^{-1},
 \quad e_{B,N}=(\Lambda_kv_+)^*Q_N\Lambda_kv_+ .\tag{OB2}
\]
Both minima include every polynomial relation and kernel lift. In
particular, $e_{B,N}$ is the minimum of $\|P\|_{m_k}^2$ over
$\deg P\le N$, $\Lambda_k[P]=\Lambda_kv_+$, and
$e_{B,N}=E_N(1-\theta_N)$ is the exact FW38 angle identity.

Set $j=k-8$, $q_-=(k-7)^2$, $\Delta=q-q_-=16k-48$ and $c_-=c-4$.
The lower cardinal $v_-$ uses $\chi_j$ and $j\rho$ in (OB1).
The conductor $\mathcal T_A$ has the literal coefficients (AC1),
without multiplication by $v_A!/\mu_{v_A}$, where
$v_A=\operatorname{ord}_0E_A\le80$. The exact maps are
\[
 C[P]=[\mathcal T_AP]\in E_j,\qquad
 \ker\Lambda_k\subset\ker C,\qquad Cv_+=a_{88}v_- .\tag{OB3}
\]
The first assertion descends because a lower root plus any degree-eight
shift is an upper root. The inclusion is the transpose of inclusion
of every conductor-multiple lower biform into the original invariant
observation image. For the last assertion, evaluation at every lower
root gives zero except at $j\rho$, where only the $(8,8)$ shift
reaches $k\rho$. This proves its value, including its coefficient.
The degree drops by at least $v_A$ by the exponential Taylor expansion.

\subsection{A consecutive-minimum estimate retaining the entire residual}
Use the original reference $d\sigma=|\Gamma(1/4+iy/2)|^2dy/(2\pi)$.
Put $Q_-(y)=i^{-q_-}\chi_j(c_-+iy)$, $dW_-=Q_-^2d\sigma$.
The polynomial $Q_-$ is real and even; since $j$ is odd, every root
has imaginary part of modulus at least $\delta$. It has no real zero.
Let $U_r$ be the real monic orthogonal polynomials for $W_-$,
$\nu_r=\|U_r\|_{W_-}^2$, and $a_r=\sqrt{\nu_r/\nu_{r-1}}$.
Their recurrence is $yU_r=U_{r+1}+a_r^2U_{r-1}$.

\begin{lemma}[Two bounds for every recurrence coefficient]
For $r\ge1$, with $B_-=2q_-/\delta+2+\pi/2$,
\[
 r/B_-\le a_r\le2(q_-+r).\tag{OB4}
\]
\end{lemma}
\begin{proof}
The digamma series \cite[5.7.6]{DLMF5} gives
$\Im\psi(1/4+ib)=\sum_{n\ge0}b/((n+1/4)^2+b^2)$.
The first term has modulus at most two. For $|b|>0$, the remaining
positive decreasing terms are bounded by
$\int_{1/4}^\infty |b|/(u^2+b^2)du\le\pi/2$; at zero they vanish.
Differentiating the actual Gamma factor gives
$(\log\sigma)'(y)=-\Im\psi(1/4+iy/2)$, so its modulus is at most
$2+\pi/2$. Also $|Q_-'/Q_-|\le q_-/\delta$ on the whole real line.
Thus $|W_-'/W_-|\le B_-$. Integration of
$(U_rU_{r-1}W_-)'$ has zero boundary values by exponential Gamma
decay. Orthogonality gives
$r\nu_{r-1}=-\int U_rU_{r-1}W_-'$; Cauchy--Schwarz proves the
lower bound. For the upper bound, $yU_{r-1}$ is a monic trial.
The exact Gamma recurrence has adjacent entries
$\sqrt{(n+1)(n+1/2)}$ and $\sqrt{n(n-1/2)}$, hence multiplication
on degree $n$ has norm at most $2(n+1)$, including its output degree.
Apply it to $Q_-U_{r-1}$, whose degree is $q_-+r-1$.
\end{proof}

Let $z_-=j\gamma-ij\delta$, $d_-=j\delta$, $R_-=|z_-|$ and
\[
 H_r=\min_{\deg P<r}\int\left|\frac1{z_--y}-P(y)\right|^2dW_-(y),
 \quad r\ge0.\tag{OB5}
\]
The empty minimum at $r=0$ means no polynomial is subtracted.
\begin{lemma}[Consecutive full residual]
\[
 \frac{a_{r+1}^2}{R_-^2+a_{r+1}^2+a_{r+2}^2}
 \le H_{r+1}/H_r\le
 \frac{a_{r+1}^2}{d_-^2+a_{r+1}^2}.\tag{OB6}
\]
\end{lemma}
\begin{proof}
The measure $W_-$ has an exponential moment near zero. Polynomials
are complete: an orthogonal vector times $W_-$ has an analytic
bilateral transform in a strip by Cauchy--Schwarz; all derivatives
at zero vanish; the identity theorem and Fourier uniqueness give zero.
In the orthonormal polynomial basis, multiplication by $y$ is a
self-adjoint Jacobi operator $J$. The coefficients of $(z_--y)^{-1}$
are $\sqrt{\nu_0}w$, $w=(z_--J)^{-1}e_0$. Therefore
$H_r=\nu_0\sum_{n\ge r}|w_n|^2$.

Remove the single edge of size $a_{r+1}$ between indices $r,r+1$.
This bounded self-adjoint finite-rank perturbation splits off a finite
prefix and a self-adjoint tail $J^{(r+1)}$. Its resolvent equation is
$w_{>r}=a_{r+1}w_r(z_--J^{(r+1)})^{-1}e_0$.
The coefficient $w_r$ cannot vanish: then this tail vanishes, and
successive recurrence equations force $w_0=0$ and contradict the
zeroth equation $(z_--J)w=e_0$.
If $L=\|(z_--J^{(r+1)})^{-1}e_0\|^2$, it follows exactly that
$H_{r+1}/H_r=a_{r+1}^2L/(1+a_{r+1}^2L)$.
The spectral probability of its unit vector $e_0$ has first moment
zero and second moment $a_{r+2}^2$. Jensen applied to the positive
variable $|z_--y|^2$ gives $L\ge(R_-^2+a_{r+2}^2)^{-1}$;
the self-adjoint resolvent estimate gives $L\le d_-^{-2}$.
Monotonicity of $x/(1+x)$ proves (OB6). The original mass $\nu_0$
has cancelled between the two norms, not been replaced by one.
\end{proof}

Write $\mathcal E_-^\sigma(L)$ for the attained norm of $v_-$ at
degree $L$ and centre $c_-$. Since
$\chi_j'(j\rho)=i^{q_--1}Q_-'(z_-)$,
\[
 v_-(c_-+iy)=\frac{Q_-(y)}{(y-z_-)Q_-'(z_-)},\quad
 \mathcal E_-^\sigma(q_--1+r)=H_r/|Q_-'(z_-)|^2.\tag{OB7}
\]
The free relation polynomials have degree below $r$, which verifies
this is the whole affine minimum. For $q-1\le N\le2q$, put
\[
 g_{k,N}=\{1+B_-^2[R_-^2+4(N-v_A+3)^2]\}^{-1}.\tag{OB8}
\]
In the $\Delta-v_A$ steps between $N-\Delta$ and $N-v_A$, (OB4)
and (OB6) bound each ratio below by (OB8): use $r+1\ge1$ and
$q_-+r+2\le N-v_A+3$. The initial degree is at least $q_--1$.
Therefore
\[
 \mathcal E_-^\sigma(N-v_A)\ge
 g_{k,N}^{\Delta-v_A}\mathcal E_-^\sigma(N-\Delta).\tag{OB9}
\]

\subsection{Exact reverse lift and finite observed lower bound}
Let $\lambda=k\rho$ and
$D_k(S)=\chi_k(S)/\chi_j(S-8\rho)$. This is a monic polynomial:
translation identifies the lower grid with the upper subgrid
$8\le a,b\le k$. Its degree is $\Delta$ and
\[
 v_+(S)=\frac{D_k(S)}{D_k(\lambda)}v_-(S-8\rho),\quad
 D_k(\lambda)=\prod_{\substack{0\le r,s\le k\\(r,s)\notin[0,k-8]^2}}
 (2r\delta+2is\gamma).\tag{OB10}
\]
These follow by differentiating the exact polynomial factorization at
$\lambda$. Every factor in the product has modulus at least
$2(k-7)\delta$. For each lower representative $P_-$ of degree
$N-\Delta$, $D_k(S)P_-(S-8\rho)/D_k(\lambda)$ is an upper
representative of degree $N$, including all of its relation terms.

For a Gamma polynomial of degree at most $n$ one has the explicit bound
\[
 \|P(\cdot+z)\|_\sigma\le\mathcal T_n(z)\|P\|_\sigma,
 \qquad\mathcal T_n(z)=e(n+1)(2n+1)^{(|z|+1)/2}.\tag{OB11}
\]
Here is a finite proof. Its generating function is
$(1+t^2)^{-1/4}e^{y\arctan t}$, by the Gamma recurrence.
Translation multiplies it by $e^{z\arctan t}$. On $|t|=r<1$,
$|\arctan t|\le\tfrac12\log((1+r)/(1-r))$.
For $n\ge1$, use $r=n/(n+1)$ in Cauchy's coefficient formula to
bound each of the first $n+1$ multiplier coefficients by
$e(2n+1)^{|z|/2}$. The weighted upper diagonal at gap $l$ has norm
at most $\sqrt{2l+1}$, from $\rho_{i+l}/\rho_i\le\sqrt{2l+1}$
with $\rho_n^2=n!/(1/2)_n$. Summing the $n+1$ diagonal norms gives
(OB11). At $n=0$ translation fixes the constant and $\mathcal T_0=e$.
The mass $\sqrt{2\pi}$ appears on both sides of the same norm.

The displacement of the reverse lift is $z_0=-8\gamma+8i\delta$.
Each root of $D_k(c+iy)$ has modulus at most
$R=k\sqrt{\delta^2+\gamma^2}$ in the $y$ coordinate.
Successive Gamma multiplication bounds, with all degrees at most $N$,
prove
\[
 \mathcal E_+^\sigma(N)\le
 \frac{\mathcal T_{N-\Delta}(z_0)^2[2(N+1)+R]^{2\Delta}}
 {|D_k(\lambda)|^2}\mathcal E_-^\sigma(N-\Delta).\tag{OB12}
\]

Use CAI16--18 at degree $2q+2$, setting
$\ell_k=a_k/[1+4(2q+2+M)^2]^M$, $u_k=A_k$, $M=25$ on this
simple-root domain. They are the original positive constants and give
$\ell_k\|P\|_\sigma^2\le\|P\|_{m_k}^2\le u_k\|P\|_\sigma^2$
uniformly on the window; $\log(u_k/\ell_k)=O_h(k+\log q)$.
In the two original centres the conductor bound is
\[
 \|\mathcal T_AP\|_{\sigma,c_-}\le\mathfrak C_{A,N}\|P\|_{\sigma,c},
 \quad\mathfrak C_{A,N}=\sum_{u,w=0}^8|a_{uw}|
 \mathcal T_N((2w-8)\gamma-i(2u-8)\delta).\tag{OB13}
\]
This follows by applying (OB11) to each original summand. Every
observation lift in (OB2) must have conductor class $a_{88}v_-$
by (OB3). Its Gamma norm is therefore at least
$|a_{88}|^2\mathcal E_-^\sigma(N-v_A)/\mathfrak C_{A,N}^2$.
Take the original full minimum and combine with (OB9)--(OB12) to obtain
\[
 1-\theta_N\ge\mathfrak l_{k,N}:=
 \frac{\ell_k}{u_k}
 \frac{|a_{88}|^2|D_k(\lambda)|^2g_{k,N}^{\Delta-v_A}}
 {\mathfrak C_{A,N}^2\mathcal T_{N-\Delta}(z_0)^2
 [2(N+1)+R]^{2\Delta}}>0.\tag{OB14}
\]
The positivity holds on the original large-period domain proved in OC;
no corner coefficient is assigned by hand. On that domain all constants
depending on $h,\varpi$ are fixed before $k$ grows. Since
$\Delta=O(k)$, $-\log g_{k,N}=O_h(\log q)$ and the translations
have fixed exponents, (OB14) proves
\[
 0\le-\log(1-\theta_N)\le-\log\mathfrak l_{k,N}
 =O_{h,\varpi}(k\log q),\qquad q-1\le N\le2q.\tag{OB15}
\]
In particular $\mathfrak l_{k,N}\le1$ follows from its proved comparison
with the actual energy ratio; it need not be separately imposed.

\subsection{Receiving the estimate and strengthening arithmetic comparison}
Insert (OB15) in the proved full-source FW35--36. With
$B_N=-\log\mathfrak l_{k,N}$, the four-angle logarithm lies in
$[-B_{q-1}-B_q,B_{2q-1}+B_{2q}]$. Thus
\[
 \log\frac{e_{B,q-1}e_{B,q}}{e_{B,2q-1}e_{B,2q}}
 =C_\partial q+O_{h,\varpi}(k\log q).\tag{OB16}
\]
Each single window has half the leading coefficient. Since
$k\log q=o(q/\log q)$, the angle contributes at neither displayed
scale. For fixed $0<s\le1$, the same substitution in FW43 gives
\[
 -\log\frac{e_{B,q-1+\lfloor sq\rfloor}}{e_{B,q-1}}
 =q\mathcal J(s)+O_{h,\varpi,s}(k\log q),\quad
 \mathcal J(s)=\int_0^sf(t)dt.\tag{OB17}
\]
Here the unchanged explicit FW44 profile is
$f(t)=\log((1+r_t)/(1-r_t))$,
$r_t=\sqrt{1-\kappa_t^2}$,
$E(\kappa_t)/(r_tK(\kappa_t))=1+t$; its logarithmic singularity
at zero is integrable, and $\mathcal J(0)=0$.
Here the complete elliptic integrals retain their usual modulus
convention
\[
 K(\kappa)=\int_0^{\pi/2}\frac{d\alpha}{\sqrt{1-\kappa^2\sin^2\alpha}},
 \qquad E(\kappa)=\int_0^{\pi/2}\sqrt{1-\kappa^2\sin^2\alpha}\,d\alpha.
\]
The original endpoint solution has $0<\kappa_t<1$ for $t>0$.
Equivalently the fully specified constant in (OB16) is
$C_\partial=2\mathcal J(1)=2\int_0^1f(t)dt$, by the two
single-window identities in FW35 and FW43. This is the same original
boundary constant, with no change of modulus or source measure.

There is a sharper arithmetic comparison than the intake's equation
(31), valid for every fixed nonzero observed class, without a corner
restriction. Minimize the two whole-polynomial CAI bounds on exactly
the same observation fibre. At each original cutoff this gives
\[
 \ell_ke_{B,N}^\sigma\le e_{B,N}^{\rm ar}\le u_ke_{B,N}^\sigma.
 \tag{OB18}
\]
For the four-cutoff logarithmic returns $\mathcal E_B^{\rm ar}$ and
$\mathcal E_B^\sigma$, sum the two numerator inequalities and subtract
the two denominator inequalities to get
\[
 |\mathcal E_B^{\rm ar}-\mathcal E_B^\sigma|
 \le2\log(u_k/\ell_k)=O_h(k+\log q).\tag{OB19}
\]
The common source mass cancels only in this equal-sign-count expression.
For a single prefix ratio the bound is $\log(u_k/\ell_k)$. This
strengthens the incoming $O(k\log q)$ comparison and is uniform over
all nonzero observation classes, even where (OB14) is not invoked.

Finally set $b_r^B=1-e_{B,N+1}/e_{B,N}$, $r=N+1-q$. The original
rank-one update FW40 identifies it, on its defined quotient domain, as
$T_r|1-V_N|^2/[d_{K,N}(1-\theta_N)]$; the undivided energy formula
defines it also when a displayed response denominator vanishes.
Positive nested minima give $0\le b_r^B<1$. Telescoping proves
\[
 -\sum_{r=0}^{q}w_r\log(1-b_r^B)
 =C_\partial q+O_{h,\varpi}(k\log q),\quad
 w_0=w_q=1,\quad w_r=2\ (0<r<q).\tag{OB20}
\]
Moreover $q^{-1}\sum_{r=0}^{q-1}-\log(1-b_r^B)\delta_{r/q}$
converges weakly on $[0,1]$ to $f(s)ds$. Its prefix sums are (OB17)
up to one boundary atom; sandwich that atom between fixed nearby
prefixes and use continuity of $\mathcal J$. The zero endpoint follows
by first using a fixed positive prefix then tending to zero; the total
mass follows from (OB16)'s single-window version. Step approximations
then prove convergence against every continuous test function.
This proves the correlated observed loss and its profile. It does not
fix the consecutive complex phases in HG7, the separate kernel
determinant, or the independent proper-source target minimum.
